\section*{第 \ref{ch:b2:fragment-orbitals} 章——片段轨道与多原子分子}

\begin{solution}{exo:b2:fragment-orbitals:1}
$h_1 + h_2$ 与 $h_1 - h_2$。平面 $xz$ 平分键角，并
垂直于分子平面，交换两个氢原子：$h_1 - h_2$ 变号，
$h_1 + h_2$ 不变。
\end{solution}

\begin{solution}{exo:b2:fragment-orbitals:2}
价层轨道六个（氧的 2s 与三个 2p，两个氢的 1s），价
电子八个，充满的轨道四个。
\end{solution}

\begin{solution}{exo:b2:fragment-orbitals:3}
充满的价层能级是一个 $a_1$ 轨道和一组三重简并的 $t_2$
轨道：两个能量，两个能带。$t_2$ 能带对应三个轨道、六个电子，
而 $a_1$ 则为一个轨道、两个电子。
\end{solution}

\begin{solution}{exo:b2:fragment-orbitals:4}
$\pi$ 键要求两个碳的 p 轨道平行，这使
两个 \ce{CH2} 平面位于同一平面内。
\end{solution}

\begin{solution}{exo:b2:fragment-orbitals:5}
氨失去的是氮上孤电子对中的一个电子，该轨道主要为\omterm{def:b2:diatomic-mos:bonding}{非键轨道}；
甲烷则来自成键的 $t_2$ 轨道。非键电子所处能级较高，
更接近原子能级，因而比成键电子更易被移去。
\end{solution}

\begin{solution}{exo:b2:fragment-orbitals:6}
线性水分子有两个简并的充满非键 p 轨道；弯曲使其中一个
降低，同时使一个成键能级略微升高：八个电子的净效果是能量降低。\ce{BeH2}
只有四个电子，充满两个\omterm{def:b2:diatomic-mos:bonding}{成键轨道}；弯曲会降低的轨道为空，
而成键能级会失去重叠：线性构型最为有利。
\end{solution}

\begin{solution}{exo:b2:fragment-orbitals:7}
与阳离子中心相邻的碳上，每个 C--H 或 C--C 键都可以与
空的 p 轨道取向一致，并带来约 $2\beta^2/\Delta$ 的稳定化。\ce{CH3+} 没有这样的
近邻；乙基、异丙基和叔丁基阳离子分别有一、二、三个甲基
与中心相邻，各提供一个取向一致的键：稳定性
依此顺序增大。
\end{solution}

\begin{solution}{exo:b2:fragment-orbitals:8}
两个能级相等时，分裂为 $2|\beta|$，两个 $\pi$ 电子共获得 $2|\beta|$，
与 $\cos\theta$ 成正比。在 $\theta = 60^\circ$ 时降为一半，
在 $90^\circ$ 时降为零。
\end{solution}

\begin{solution}{exo:b2:fragment-orbitals:9}
硼烷有六个价电子，它们充满三个\omterm{def:b2:diatomic-mos:bonding}{成键轨道}，
垂直于平面的硼 p 轨道保持为空，锥化没有任何收益。
氨多两个电子，占据平面构型中的\omterm{def:b2:diatomic-mos:bonding}{非键轨道}；
与水中一样，锥化使其与 2s 轨道及氢的组合相混合，
能量随之降低。
\end{solution}

\begin{solution}{exo:b2:fragment-orbitals:10}
当 $c_1 = c_2 = c_3$ 时，久期方程给出 $\alpha + 2\beta$；与之正交的两个组合
给出简并的 $\alpha - \beta$。两个电子占据 $\alpha + 2\beta$，能量为
$2\alpha + 4\beta$，低于 $\ce{H2} + \ce{H+}$（$2\alpha + 2\beta$），因为 $\beta < 0$。
\end{solution}

\begin{solution}{exo:b2:fragment-orbitals:11}
$\pi_1$：三个系数同号，无节点（成键）；$\pi_2$：
两端碳上符号相反，中心有节点（非键）；$\pi_3$：符号交替，
两个节点（反键）。该阴离子有四个 $\pi$ 电子，$\pi_2$ 为其
最高占据能级，电子密度只分布在两端碳上。
\end{solution}

\begin{solution}{exo:b2:fragment-orbitals:12}
在包含该轴的两个平面中，各有三个平行的 p 轨道（O、C、O），
分别组合成一个\omterm{def:b2:diatomic-mos:bonding}{成键轨道}、一个\omterm{def:b2:diatomic-mos:bonding}{非键轨道}（仅在氧上）和一个\omterm{def:b2:diatomic-mos:bonding}{反键轨道}：共六个
$\pi$ 轨道。两个\omterm{def:b2:diatomic-mos:bonding}{成键轨道}与两个\omterm{def:b2:diatomic-mos:bonding}{非键轨道}均已充满：八个
$\pi$ 电子（在路易斯结构中对应两个 $\pi$ 键与两个氧的孤电子对）。
\end{solution}

\begin{solution}{pb:b2:fragment-orbitals:1}
\textbf{1.} 八个：氧提供六个，每个氢各提供一个。
\textbf{2.} $h_1 + h_2$ 与 $h_1 - h_2$。
\textbf{3.} 2s 与 $2\mathrm p_z$。
\textbf{4.} $2\mathrm p_y$。
\textbf{5.} $2\mathrm p_x$，垂直于分子平面。
\textbf{6.} 六个；其中四个充满。
\textbf{7.} $h_1 - h_2$，穿过氧原子处垂直于轴的
平面时变号，与 $2\mathrm p_z$ 相同。
\textbf{8.} $h_1 + h_2$。
\textbf{9.} $2\mathrm p_x$ 与 $2\mathrm p_y$，均垂直于轴，构成一对
简并的\omterm{def:b2:diatomic-mos:bonding}{非键轨道}。
\textbf{10.} （2s，$h_1 + h_2$）成键$^2$，（$2\mathrm p_z$，$h_1 - h_2$）成键$^2$，然后
$2\mathrm p_x^2\,2\mathrm p_y^2$。
\textbf{11.} 否：简并的一对已充满，所有电子都已成对。
\textbf{12.} 位于氧 2p 轨道的能量处，它是\omterm{def:b2:diatomic-mos:bonding}{非键轨道}。
\textbf{13.} $2\mathrm p_y$（线性分子的轴沿 $z$，弯曲在 $yz$ 平面内）。
\textbf{14.} 它与 2s 及 $h_1 + h_2$ 混合并降低：即 $3a_1$ 能级。
\textbf{15.} （$2\mathrm p_z$，$h_1 - h_2$）成键能级，它随氢原子离开轴线
而损失部分重叠，成为 $1b_2$。
\textbf{16.} 降低的充满能级所获得的能量多于另一能级所损失的：八个
电子时，弯曲构型更稳定。
\textbf{17.} $104.48^\circ$，介于 $90^\circ$（键仅由 p 轨道构成）与 $109.5^\circ$ 之间：
2s 轨道参与了成键。
\textbf{18.} $2\mathrm p_x$，仍为非键的 $1b_1$。
\textbf{19.} 四个。
\textbf{20.} 来自 $1b_1$，即氧的非键 $2\mathrm p_x$ 轨道。
\textbf{21.} 移去一个非键电子几乎不改变成键：离子的几何构型与分子相同，
被激发的振动能很少。移去
一个成键电子则改变键长和键角，谱带
分布于众多振动能级上。
\textbf{22.} \qty{12.62}{eV} 对比 \qty{13.62}{eV}：水中的氧带有一部分负电荷，
由氢原子拉来，使其 2p 能级升高。
\textbf{23.} 它们并非两个等能级：其两种组合分别为 $3a_1$ 和
$1b_1$ 轨道，能量不同。两种图像描述的是同一总
电子密度；只有离域轨道才给出电离能。
\textbf{24.} 水有\textbf{四}个价层电离谱带，最低者位于
$\boldsymbol{\approx \qty{12.6}{eV}}$。
\end{solution}
